\section*{Chapitre \ref{ch:g10:refraction} --- Réfraction de la lumière}

\begin{solution}{exo:g10:refraction:1}
Les angles se mesurent depuis la \omterm{def:g10:refraction:angles}{normale}~: l'\omterm{def:g10:refraction:angles}{angle d'incidence}
est $90^\circ - 25^\circ = 65^\circ$, donc l'angle de réflexion est
aussi $65^\circ$. Les faisceaux incident et réfléchi font un angle de
$2 \times 65^\circ = 130^\circ$ l'un avec l'autre.
\end{solution}

\begin{solution}{exo:g10:refraction:2}
\emph{1.} $v = c/n$~: dans l'eau
$v = \qty{3.00e8}{m/s}/1.33 = \qty{2.26e8}{m/s}$~; dans le diamant
$v = \qty{3.00e8}{m/s}/2.42 = \qty{1.24e8}{m/s}$.

\emph{2.} $n = \dfrac{c}{v} = \dfrac{\qty{3.00e8}{m/s}}
{\qty{1.97e8}{m/s}} = 1.52$~: verre.
\end{solution}

\begin{solution}{exo:g10:refraction:3}
$\sin i_2 = \dfrac{\sin 40^\circ}{1.50} = \dfrac{0.643}{1.50} = 0.429$,
donc $i_2 = 25.4^\circ$. Par réversibilité des trajets lumineux, un
rayon à l'intérieur du verre frappant la surface à $25.4^\circ$ sort
dans l'air exactement à $40^\circ$.
\end{solution}

\begin{solution}{exo:g10:refraction:4}
$\sin i_2 = \dfrac{\sin 55^\circ}{1.33} = \dfrac{0.819}{1.33} = 0.616$,
donc $i_2 = 38.0^\circ$. Le rayon se courbe \emph{vers} la
\omterm{def:g10:refraction:angles}{normale}, parce que l'eau est le milieu plus lent ($n_2 > n_1$).
\end{solution}

\begin{solution}{exo:g10:refraction:5}
$n = \dfrac{\sin 45^\circ}{\sin 32^\circ} = \dfrac{0.707}{0.530}
= 1.33$~: le liquide est de l'eau.
\end{solution}

\begin{solution}{exo:g10:refraction:6}
\emph{1.} Les cinq rapports sont
\[
\frac{0.342}{0.233} = 1.47,\quad
\frac{0.500}{0.334} = 1.50,\quad
\frac{0.643}{0.431} = 1.49,\quad
\frac{0.766}{0.515} = 1.49,\quad
\frac{0.866}{0.581} = 1.49 :
\]
constants à l'erreur de mesure près --- \omterm{thm:g10:refraction:snell}{loi de Snell--Descartes}.

\emph{2.} $n \approx 1.49$~: plexiglas.

\emph{3.} La \omterm{thm:g10:refraction:snell}{loi de Snell--Descartes} prédit
$\sin i_2 = \sin i_1 / n$~: sinus contre sinus, les données doivent
s'aligner sur une droite passant par l'origine, qu'une règle vérifie
d'un coup d'œil. Tracés angle contre angle, les données se trouvent
sur une courbe~; aux petits angles la courbe est trompeusement proche
d'une droite, ce qui a précisément trompé Ptolémée.
\end{solution}

\begin{solution}{exo:g10:refraction:7}
$\sin i_c = n_2/n_1$ avec $n_2 = 1.00$~:
eau $\sin i_c = 1/1.33 = 0.752$, $i_c = 48.8^\circ$~;
verre $\sin i_c = 1/1.50 = 0.667$, $i_c = 41.8^\circ$~;
diamant $\sin i_c = 1/2.42 = 0.413$, $i_c = 24.4^\circ$.

De l'air vers l'eau, $\sin i_2 = \sin i_1/1.33 \leq 1/1.33 < 1$ pour
toute incidence~: un angle réfracté existe toujours, donc rien de
spécial n'arrive jamais --- la \omterm{def:g10:refraction:critical}{réflexion totale interne} exige
$n_1 > n_2$.
\end{solution}

\begin{solution}{exo:g10:refraction:8}
\emph{1.} $d' = \dfrac{d}{n} = \dfrac{\qty{2.0}{m}}{1.33}
= \qty{1.5}{m}$.

\emph{2.} Le poisson apparaît à $\qty{40}{cm}/1.33 = \qty{30}{cm}$~:
l'image est \emph{au-dessus} du poisson, donc le héron doit frapper
\emph{en dessous} de l'image qu'il voit. Les hérons, non encombrés de
la \omterm{thm:g10:refraction:snell}{loi de Snell--Descartes}, apprennent cela par essais et
erreurs.
\end{solution}

\begin{solution}{exo:g10:refraction:9}
\emph{1.} $\sin r = \dfrac{\sin 50^\circ}{1.50} = \dfrac{0.766}{1.50}
= 0.511$, donc $r = 30.7^\circ$.

\emph{2.} Verre vers air à $30.7^\circ$~:
$\sin i_2 = 1.50 \times 0.511 = 0.766$, donc $i_2 = 50^\circ$ ---
l'angle de sortie égale l'angle d'entrée.

\emph{3.} Les deux \omterm{def:g10:refraction:refraction}{réfractions} s'annulent~: le rayon sort
\emph{parallèle} à sa direction d'origine, seulement décalé
latéralement. C'est pourquoi une vitre ne distord pas la scène, elle
ne fait que (imperceptiblement) la translater.
\end{solution}

\begin{solution}{exo:g10:refraction:10}
$\sin 55^\circ = 0.819$. Rouge~:
$\sin i_2 = 0.819/1.61 = 0.509$, $i_2 = 30.6^\circ$. Bleu~:
$\sin i_2 = 0.819/1.66 = 0.493$, $i_2 = 29.6^\circ$. Les deux couleurs
sont séparées de $1.0^\circ$~; le bleu, avec le plus grand indice, se
courbe davantage et finit plus près de la \omterm{def:g10:refraction:angles}{normale}.
\end{solution}

\begin{solution}{exo:g10:refraction:11}
\emph{1.} $\sin i_c = \dfrac{1.46}{1.48} = 0.986$, donc
$i_c = 80.6^\circ$.

\emph{2.} Un rayon à $i_c$ de la \omterm{def:g10:refraction:angles}{normale} fait $90^\circ - i_c$ avec
l'axe, donc son trajet le long d'une longueur d'axe $L$ est
$\dfrac{L}{\sin i_c}$~:
\[
\qty{1000}{m} \times \frac{1.48}{1.46} = \qty{1013.7}{m},
\]
environ \qty{14}{m} de plus que le rayon axial.
\end{solution}

\begin{solution}{exo:g10:refraction:12}
\emph{1.} La lumière rasant depuis l'horizon arrive à $i_1$ proche de
$90^\circ$~; elle se réfracte à $\sin i_2 = \sin 90^\circ / 1.33 = 0.752$,
c'est-à-dire $i_2 = 48.8^\circ$ --- l'\omterm{def:g10:refraction:critical}{angle critique}, comme la
réversibilité l'exige. Les rayons de tout le ciel atteignent donc l'œil
dans un cône de demi-angle $48.8^\circ$.

\emph{2.} Le bord du disque est là où ce cône rencontre la surface~:
\[
r = d \tan i_c = \qty{2.0}{m} \times \tan 48.8^\circ
= \qty{2.0}{m} \times 1.14 = \qty{2.3}{m}.
\]

\emph{3.} À l'extérieur du disque, la lumière d'en dessous frappe la
surface au-delà de l'\omterm{def:g10:refraction:critical}{angle critique} et est totalement réfléchie~: le
plongeur voit un miroir argenté montrant le fond de la piscine.
\end{solution}

\begin{solution}{exo:g10:refraction:13}
\emph{1.} Rouge~: $\sin i_2 = 1.51 \times \sin 30^\circ = 0.755$, donc
$i_2 = 49.0^\circ$. Bleu~: $\sin i_2 = 1.53 \times 0.500 = 0.765$,
donc $i_2 = 49.9^\circ$. Le faisceau blanc s'éventaille sur
$49.9^\circ - 49.0^\circ = 0.9^\circ$ --- un \omterm{def:g10:light-spectra:white}{spectre} de prisme
issu d'une seule \omterm{def:g10:refraction:refraction}{réfraction}.

\emph{2.} À $45^\circ$~: $\sin i_2 = 1.51 \times 0.707 = 1.07 > 1$
(et $1.53 \times 0.707 = 1.08$ pour le bleu). Aucun angle réfracté
n'existe pour aucune des deux couleurs~: l'\omterm{def:g10:refraction:critical}{angle critique} est
$\arcsin(1/1.51) = 41.5^\circ < 45^\circ$, donc la grande face réfléchit
totalement le faisceau entier.
\end{solution}

\begin{solution}{exo:g10:refraction:14}
\emph{1.} \omterm{def:g10:refraction:critical}{Angle critique} verre--air~: $\sin i_c = 1/1.50$, donc
$i_c = 41.8^\circ < 45^\circ$~: le faisceau est totalement réfléchi ---
$100\%$ de la lumière, sans revêtement à ternir. Un miroir métallique
perd quelques pourcents à chaque rebond, et un périscope en a besoin
de deux.

\emph{2.} Verre--eau~: $\sin i_c = \dfrac{1.33}{1.50} = 0.887$, donc
$i_c = 62.5^\circ$. Maintenant $45^\circ < i_c$~: l'incidence est sous
le critique, la plupart de la lumière se réfracte dans l'eau, et le
miroir --- donc l'image --- échoue.

\emph{3.} $\sin i_2 = \dfrac{1.50}{1.33} \times \sin 45^\circ
= 1.128 \times 0.707 = 0.798$, donc $i_2 = 52.9^\circ$ dans l'eau.
\end{solution}

\begin{solution}{exo:g10:refraction:15}
\emph{1.} $v = \dfrac{\qty{3.00e8}{m/s}}{1.47} = \qty{2.04e8}{m/s}$,
donc
\[
t = \frac{\qty{6.2e6}{m}}{\qty{2.04e8}{m/s}} = \qty{30.4}{ms} \approx
\qty{30}{ms}.
\]

\emph{2.} Radio~: $t = \qty{6.2e6}{m}/\qty{3.00e8}{m/s} =
\qty{20.7}{ms}$. Le verre coûte $30.4 - 20.7 = \qty{9.7}{ms}$ dans
chaque sens, environ \qty{19}{ms} par aller-retour.

\emph{3.} Abaisser l'indice --- les fibres à \omterm{def:g10:refraction:fiber}{cœur creux} guident la
lumière à travers l'air ($n \approx 1.00$), et les liaisons micro-ondes
à travers l'atmosphère font de même --- ou raccourcir le trajet~: un
câble posé plus près de la route orthodromique entre les deux villes.
Les deux se vendent effectivement, au prix de la milliseconde.
\end{solution}

\begin{solution}{pb:g10:refraction:1}
\textbf{1.} $v_{\text{water}} = \qty{3.00e8}{m/s}/1.33 =
\qty{2.26e8}{m/s}$~; $v_{\text{diamond}} = \qty{3.00e8}{m/s}/2.42 =
\qty{1.24e8}{m/s}$ --- le diamant ralentit la lumière d'un facteur
$n = 2.42$.

\textbf{2.} $\sin i_2 = \sin 40^\circ / 1.33 = 0.643/1.33 = 0.483$,
donc $i_2 = 28.9^\circ$.

\textbf{3.} $\sin i_2 = 1.33 \times \sin 28.9^\circ = 1.33 \times
0.483 = 0.643$, donc $i_2 = 40^\circ$~: la lumière reprend son
trajet --- réversibilité des rayons lumineux.

\textbf{4.} $\sin i_c = 1/1.33 = 0.752$, donc $i_c = 48.8^\circ$. À
$60^\circ > i_c$, la \omterm{thm:g10:refraction:snell}{loi de Snell--Descartes} exigerait
$\sin i_2 = 1.33 \sin 60^\circ = 1.15 > 1$~: aucun
\omterm{def:g10:refraction:refraction}{rayon réfracté}~; la surface réfléchit totalement la lumière
vers le bas.

\textbf{5.} $v = \qty{3.00e8}{m/s}/1.0003 = \qty{2.999e8}{m/s}$~: la
lumière dans l'air est plus lente que dans le vide de $0.03\%$, bien
en dessous de la précision de nos mesures d'angles, donc
$n_{\text{air}} = 1.00$ est sans conséquence.

\textbf{6.} $d' = \qty{3.0}{m}/1.33 = \qty{2.3}{m}$~: la piscine cache
\qty{70}{cm} d'elle-même.

\textbf{7.} Deux rayons quittant la pièce se réfractent à la surface,
s'écartant \emph{de} la \omterm{def:g10:refraction:angles}{normale}, donc ils divergent plus
fortement au-dessus de l'eau qu'en dessous. L'œil prolonge les deux
rayons sortants en arrière en lignes droites~; ils se croisent
\emph{au-dessus} de la pièce, et ce croisement est où l'œil place
l'image.

\textbf{8.} $r = d \tan i_c = \qty{3.0}{m} \times \tan 48.8^\circ =
\qty{3.0}{m} \times 1.14 = \qty{3.4}{m}$~: un disque d'environ
\qty{6.9}{m} de diamètre contient le ciel entier.

\textbf{9.} À $30^\circ$~: $\sin i_2 = 1.33 \times 0.500 = 0.665$, le
rayon sort à $41.7^\circ$. À $45^\circ$~: $\sin i_2 = 1.33 \times
0.707 = 0.940$, il sort à $70.1^\circ$, presque rasant. À
$50^\circ$~: $1.33 \times 0.766 = 1.02 > 1$ --- \omterm{def:g10:refraction:critical}{réflexion totale
interne}, le rayon plonge à nouveau dans la piscine.

\textbf{10.} L'image est au-dessus du vrai poisson
(\cref{prop:g10:refraction:depth}), donc viser \emph{en dessous}. Et
oui~: le poisson voit aussi le pêcheur déplacé, comprimé vers la
verticale dans la fenêtre de Snell --- chacun voit l'autre là où
l'autre n'est pas.

\textbf{11.} Diamant~: $\sin i_c = 1/2.42 = 0.413$, donc
$i_c = 24.4^\circ$, contre $41.1^\circ$ pour le verre~: le cône
d'échappement du diamant n'est même pas la moitié aussi large.

\textbf{12.} Dans le diamant, $30^\circ > 24.4^\circ$~: \omterm{def:g10:refraction:critical}{réflexion
totale interne}, le rayon reste dans la pierre. Dans le verre,
$30^\circ < 41.1^\circ$~: le rayon fuit par l'arrière. Les facettes
d'une taille brillant sont inclinées pour que la lumière entrante
frappe presque toujours au-delà de $24.4^\circ$ et rebondisse jusqu'à
sortir vers le haut vers l'œil --- une imitation en verre, avec son
large cône d'échappement, laisse simplement la lumière tomber à
travers.

\textbf{13.} $\sin 45^\circ = 0.707$. Rouge~: $0.707/2.41 = 0.293$,
donc $i_2 = 17.1^\circ$. Bleu~: $0.707/2.45 = 0.289$, donc
$i_2 = 16.8^\circ$. Écart à l'entrée~: environ $0.3^\circ$.

\textbf{14.} $\sin 17.0^\circ = 0.292$. Rouge~: $\sin i_2 = 2.41
\times 0.292 = 0.705$, donc $i_2 = 44.8^\circ$. Bleu~:
$\sin i_2 = 2.45 \times 0.292 = 0.716$, donc $i_2 = 45.8^\circ$.
L'éventail a maintenant $1.0^\circ$ de large~: la \omterm{def:g10:refraction:refraction}{réfraction}
de sortie, travaillant près de l'\omterm{def:g10:refraction:critical}{angle critique}, a étiré
l'écart d'entrée de $0.3^\circ$ d'un facteur trois. Ce jet de couleurs
étiré est le \emph{feu} du diamant.

\textbf{15.} Diamant--eau~: $\sin i_c = 1.33/2.42 = 0.550$, donc
$i_c = 33.3^\circ$. Le rayon à $30^\circ$, piégé lorsque la facette
faisait face à l'air, s'échappe maintenant ($30^\circ < 33.3^\circ$)
dans le film d'eau~: chaque goutte, empreinte ou film de graisse élève
l'indice extérieur, élargit le cône d'échappement et vide l'éclat ---
d'où le chiffon du joaillier.

\textbf{16.} $\sin i_c = \dfrac{1.46}{1.48} = 0.986$, donc
$i_c = 80.6^\circ$.

\textbf{17.} La \omterm{def:g10:refraction:critical}{réflexion totale interne} a lieu \emph{à la
surface}, donc la surface doit rester parfaite~: sur un fil nu, chaque
grain de poussière, rayure, goutte d'eau ou doigt qui touche élève
l'indice extérieur en ce point et saigne la lumière. La
\omterm{def:g10:refraction:fiber}{gaine} enterre le miroir dans du verre propre où rien ne peut
l'atteindre --- et, en prime, protège mécaniquement le
\omterm{def:g10:refraction:fiber}{cœur} fin comme un cheveu.

\textbf{18.} $v = \qty{3.00e8}{m/s}/1.48 = \qty{2.03e8}{m/s}$, donc
\[
t = \frac{\qty{6.0e6}{m}}{\qty{2.03e8}{m/s}} = \qty{29.6}{ms} \approx
\qty{30}{ms}.
\]

\textbf{19.} Facteur de trajet $1/\sin i_c = 1.48/1.46 = 1.0137$~: au
pire $\qty{6000}{km} \times 0.0137 \approx \qty{82}{km}$ de verre
supplémentaire, coûtant $\qty{8.2e4}{m} / \qty{2.03e8}{m/s} \approx
\qty{0.4}{ms}$ --- le zigzag est presque gratuit.

\textbf{20.} Ping~: $2 \times (29.6 + 0.4) \approx \qty{60}{ms}$.
Débit~: $\dfrac{10^{13}}{2.4 \times 10^8} \approx 40\,000$ copies de
ce livre par seconde. En une phrase~: en piégeant la lumière dans le
verre par \omterm{def:g10:refraction:critical}{réflexion totale interne}, la
\omterm{thm:g10:refraction:snell}{loi de Snell--Descartes} transporte une bibliothèque à
travers l'océan chaque seconde, soixante millisecondes aller-retour.
\end{solution}
